\section{Riccati 型の四次元 AS-regular 射影模型}
\label{sec:riccati-p3-link}

Riccati 成分には、通常の $\PP^3$ を古典ファイバーとする
AS-regular コンパクト化も構成できる。
ここで得られる $\PP^3$ 上の Poisson 構造は、
平面の次数2葉層の線形引き戻し成分 $S(2,3)$ に属する特別な族である。

\begin{proposition}\label{prop:riccati-p3-quantum-link}
Riccati 成分の一般元に対して、その微分型 Ore 量子化と
非可換双有理同値な四次元 quadratic AS-regular 代数の
一変数代数族が存在する。
この族の Hilbert 級数は $(1-s)^{-4}$、
古典的なファイバーは $\CC[Z,X,T,V]$ である。
分数体の同型の古典極限は、Riccati 成分から
$S(2,3)$ 内の特別な族への Poisson 双有理変換を与える。
\end{proposition}
\begin{proof}
元の座標を $(x,t,y)$ とし、
\[
 \partial=f(x)\partial_x+(h_0(t)+xh_1(t))\partial_t,
 \qquad f(x)=a(x-r_1)(x-r_2)(x-r_3)
\]
と書く。三つの根は相異なるとし、$d=r_3-r_1$、$k=x-r_3$ と置く。
有理座標
\[
 \xi=\frac{x-r_1}{x-r_3},\qquad \eta=\frac{y}{x-r_3}
\]
に移ると、Poisson 構造は
\[
 \left(F(\xi)\partial_\xi
     +(H_0(t)+\xi H_1(t))\partial_t\right)\wedge\partial_\eta
\]
となる。ただし
\begin{align*}
 F(\xi)&=-a\xi\bigl((r_3-r_2)\xi+r_2-r_1\bigr),\\
 H_0(t)&=-\frac{h_0(t)+r_1h_1(t)}d,
 &H_1(t)&=\frac{h_0(t)+r_3h_1(t)}d.
\end{align*}
実際、$\{\xi,\eta\}=f(x)\xi'(x)/k$、
$\{t,\eta\}=(h_0+xh_1)/k$ であり、
$1/k=(\xi-1)/d$ を代入すればこれらの式を得る。

一般の $h_0,h_1$ に対して $H_1(t)\partial_t$ は半単純である。
これは $H_1$ の二つの零点が異なるという Zariski 開条件であり、
$h_0+r_3h_1$ の係数は任意の二次式の係数を動けるので非空である。
その二つの零点を $s_1,s_2$ として
$H_1(t)=b(t-s_1)(t-s_2)$ と書き、
$\tau=(t-s_1)/(t-s_2)$ に移ると
\[
 H_1(t)\partial_t=\kappa\tau\partial_\tau,
 \qquad \kappa=b(s_1-s_2).
\]
$H_0(t)\partial_t$ は $\PP^1$ の大域ベクトル場なので、
同じ座標では二次以下の式 $A(\tau)\partial_\tau$ になる。
従って、座標を再び $(x,t,y)$ と書き直すと、
\begin{equation}\label{eq:riccati-quadratic-normal}
 \pi=\left(F(x)\partial_x+H(x,t)\partial_t\right)\wedge\partial_y,
 \qquad H(x,t)=A(t)+\kappa xt,
\end{equation}
を得る。$F,H$ は全次数が高々2である。
これらの正規化を族で行うときは、根を標識する有限 \'etale 被覆を用いる。

量子座標の変換も厳密に定義できる。
元の微分型 Ore 分数体では、左除算を用いて
\[
 \xi=\frac{x-r_1}{x-r_3},\qquad
 \eta=(x-r_3)^{-1}y
\]
と置く。係数体は可換なので、任意の $a\in\CC(x,t)$ に対して
\[
 [a,\eta]=\hbar k^{-1}\partial(a)
\]
である。続く $t$ の M\"obius 変換も係数体の同型なので、
\eqref{eq:riccati-quadratic-normal} の導分を用いた
Ore 分数体と元の分数体は $\CC(\hbar)$ 上で同型である。
逆写像は $x=(r_3\xi-r_1)/(\xi-1)$、$y=k\eta$ と、
$t$ の逆 M\"obius 変換によって与えられる。

$F,H$ の二次斉次化を
\[
 F^{\mathrm h}=Z^2F(X/Z),\qquad
 H^{\mathrm h}=Z^2A(T/Z)+\kappa XT
\]
とし、
\[
 \ell=-\frac{\partial_XF^{\mathrm h}+\partial_TH^{\mathrm h}}4,
 \qquad
 \Delta=(F^{\mathrm h}+X\ell)\partial_X
          +(H^{\mathrm h}+T\ell)\partial_T+Z\ell\partial_Z
\]
と置く。$\ell$ は一次式なので、Euler の恒等式から
\[
 \operatorname{div}(\Delta)
  =\partial_XF^{\mathrm h}+\partial_TH^{\mathrm h}+4\ell=0.
\]
全生成元の次数を1として
\[
 \mathcal A=\CC[\hbar,Z,X,T][V;-\hbar\Delta]
\]
を定める。これは homogeneous differential Ore extension であり、
PBW 基底 $Z^iX^jT^kV^\ell$ を持つ。
従って次数ごとに自由で、Hilbert 級数は $(1-s)^{-4}$ である。
各ファイバーは noetherian domain であり、
Ore 拡大の定理 \cite[定理0.2、\S1.2--1.3]{LiuWangWu}
によって大域次元4の AS-regular 代数となる。

$Z$ は正規元であり、
\[
 VZ=Z(V-\hbar\ell)
\]
となる。従って $Z$ のべきを Ore 局所化できる。
局所化の次数零部分において
\[
 x=X/Z,\qquad t=T/Z,\qquad y=Z^{-1}V
\]
と取ると、
\[
 \Delta(x)=ZF(x),\qquad \Delta(t)=ZH(x,t)
\]
なので
\[
 [x,t]=0,\qquad [x,y]=\hbar F(x),\qquad
 [t,y]=\hbar H(x,t).
\]
これらは次数零部分を生成する。
PBW 基底の $V$ 次数を比較すると追加の関係はなく、
次数零部分は正規化後の Riccati Ore 代数そのものである。

最後に $S(2,3)$ への所属を確認する。
アフィンでの1形式は $H\,dx-F\,dt$ であり、同次化は
\[
 \omega=ZH^{\mathrm h}\,dX-ZF^{\mathrm h}\,dT
            +(TF^{\mathrm h}-XH^{\mathrm h})\,dZ.
\]
これは $V$ を含まず、$\PP^2_{[Z:X:T]}$ の次数2葉層の
線形射影による引き戻しである。
一般の係数について $\gcd(F^{\mathrm h},H^{\mathrm h})=1$ であり、
$F=f_2x^2+f_1x$、$A=a_2t^2+a_1t+a_0$ と書くと
\[
 (TF^{\mathrm h}-XH^{\mathrm h})|_{Z=0}
    =XT\bigl((f_2-\kappa)X-a_2T\bigr)\ne0.
\]
従って $\omega$ の係数に共通因子がなく、葉層の次数は正確に2である。
これが主張した成分内の特別な族を与える。
\end{proof}
